% Part: normal-modal-logic
% Chapter: axioms-systems
% Section: consistency

\documentclass[../../../include/open-logic-section]{subfiles}

\begin{document}

\olfileid{nml}{prf}{con}

\olsection{સુસંગતતા}

સુસંગતતા એ !!{formula}sના ગણનો અગત્યનો ગુણધર્મ છે. જો !!{formula}sના કોઈ ગણમાંથી
વિસંવાદ, જેમ કે~$\lfalse$, !!{derivable} હોય, તો તે ગણ અસુસંગત છે;
નહિ તો સુસંગત છે. ગણ અસુસંગત હોય, તો તેનાં બધાં !!{formula}s કોઈ
નિદર્શનના એક જગતમાં એકસાથે સત્ય હોઈ શકતાં નથી. પૂર્ણતાના પ્રમેય માટે
અમે વિપરીત વિધાન સાબિત કરીએ છીએ: દરેક સુસંગત ગણ કોઈ નિદર્શનના એક
જગતમાં સત્ય છે, ખાસ કરીને ``કૅનોનિકલ નિદર્શન''માં.

\begin{defn}
  ગણ $\Gamma$ એ તંત્ર~$\Sigma$ને અનુલક્ષીને \emph{સુસંગત} છે, અથવા
  આપણે કહીએ તેમ, $\Sigma$-સુસંગત છે, જો અને માત્ર જો $\Gamma
  \Proves/[\Sigma] \lfalse$.
\end{defn}

દાખલા તરીકે, ગણ $\{ \Box(p \lif q), \Box p, \lnot\Box q \}$
વિધાનાત્મક તર્કને અનુલક્ષીને સુસંગત છે, પરંતુ \Log{K}-સુસંગત નથી.
એ જ રીતે, ગણ $\{ \Diamond p, \Box\Diamond
p \lif q, \lnot q \}$ એ \Log{K5}-સુસંગત નથી.

\begin{prop}\ollabel{prop:consistencyfacts}
  માનો કે $\Gamma$ એ !!{formula}sનો ગણ છે. ત્યારે:
  \begin{enumerate}
  \item $\Gamma$ એ $\Sigma$-સુસંગત છે જો અને માત્ર જો એવું કોઈ
    !!{formula}~$!A$ હોય કે $\Gamma \Proves/[\Sigma]
    !A$.
  \item \ollabel{prop:consistencyfacts-b}%
    $\Gamma \Proves[\Sigma] !A$ જો અને માત્ર જો $\Gamma \cup \{
    \lnot!A \}$ એ $\Sigma$-સુસંગત ન હોય.
  \item \ollabel{prop:consistencyfacts-c}%
    જો $\Gamma$ એ $\Sigma$-સુસંગત હોય, તો કોઈપણ !!{formula}
    $!A$ માટે કાં $\Gamma \cup \{ !A \}$ એ $\Sigma$-સુસંગત છે,
    કાં $\Gamma \cup \{ \lnot!A \}$ એ $\Sigma$-સુસંગત છે.
  \end{enumerate}
\end{prop}

\begin{proof}
  આ હકીકતો પરંપરાગત વિધાનાત્મક તર્કથી સહેલાઈથી ફલિત થાય છે. આપણે
  \olref{prop:consistencyfacts-c}ની દલીલ આપીએ. પ્રતિપક્ષી વિધાનથી
  આગળ વધો અને માનો કે $\Gamma \cup \{ !A \}$ તથા
  $\Gamma \cup \{ \lnot!A \}$માંથી એકેય $\Sigma$-સુસંગત નથી.
  ત્યારે \olref{prop:consistencyfacts-b} મુજબ $\Gamma, !A \Proves[\Sigma]
  \lfalse$ અને $\Gamma, \lnot !A \Proves[\Sigma] \lfalse$ બંને.
  નિષ્પત્તિ પ્રમેયથી $\Gamma \Proves[\Sigma] !A \to \lfalse$ અને
  $\Gamma \Proves[\Sigma] \lnot!A \lif \lfalse$. પરંતુ $(!A \lif
  \lfalse) \lif ((\lnot!A \lif \lfalse) \lif \lfalse)$ પુનરુક્તિનો
  દૃષ્ટાંત છે; તેથી
  \olref[prp]{prop:derivabilityfacts}\olref[prp]{prop:derivabilityfacts-ruleT} મુજબ
  $\Gamma \Proves[\Sigma] \lfalse$.
\end{proof}

\end{document}
